% Compilar duas vezes, a partir desta pasta:
%   xelatex -shell-escape apostila-ecolab-opengl.tex
%   xelatex -shell-escape apostila-ecolab-opengl.tex
\documentclass[11pt,a4paper]{article}

\usepackage{mathtools}
\usepackage{fontspec}
\usepackage{unicode-math}
\setmainfont{TeX Gyre Pagella}
\setsansfont{TeX Gyre Heros}[Scale=0.94]
\setmonofont{JetBrains Mono}[Scale=0.84, HyphenChar=None]
\setmathfont{TeX Gyre Pagella Math}

\usepackage[brazilian]{babel}
\usepackage{microtype}
\usepackage{geometry}
\geometry{a4paper, margin=2.15cm, top=2.05cm, bottom=2.05cm, headheight=14.5pt}

\usepackage{xcolor}
\usepackage{tikz}
\usetikzlibrary{arrows.meta,positioning,calc}
\usepackage{booktabs}
\usepackage{tabularx}
\usepackage{array}
\usepackage{enumitem}
\usepackage{caption}
\captionsetup{font=small, labelfont={sf,bf}, skip=0.45em, hypcap=false}
\usepackage{needspace}
\usepackage{tcolorbox}
\tcbuselibrary{listings,minted,skins,breakable}
\usepackage{minted}
\usepackage{fancyhdr}
\usepackage{titlesec}
\usepackage{hyperref}
\usepackage{bookmark}

\definecolor{ink}{HTML}{1C2420}
\definecolor{muted}{HTML}{5C665E}
\definecolor{laranja}{HTML}{D67E2A}
\definecolor{laranjaesc}{HTML}{A85A18}
\definecolor{verde}{HTML}{6A8F4E}
\definecolor{verdeesc}{HTML}{3E5A2C}
\definecolor{codebg}{HTML}{F7F4EE}
\definecolor{caixafundo}{HTML}{F3F6F1}
\definecolor{contafundo}{HTML}{FBF6EE}
\definecolor{alertafundo}{HTML}{F8F1E8}
\definecolor{linhasuave}{HTML}{E4DDD2}
\definecolor{papel}{HTML}{FBF8F3}

\hypersetup{
  colorlinks=true,
  linkcolor=verdeesc,
  urlcolor=laranjaesc,
  pdftitle={A logo e-colab em OpenGL},
  pdfauthor={Apostila de computação gráfica},
  pdfkeywords={OpenGL, matrizes, petala, Phong, Gouraud, Z-buffer}
}

\setlist{itemsep=0.22em, topsep=0.35em, parsep=0.1em}
\setlist[itemize]{leftmargin=1.15em, label=\textcolor{laranja}{\textbullet}}
\setlist[enumerate]{leftmargin=1.45em, label=\textcolor{laranjaesc}{\arabic*.}}

\setlength{\parskip}{0.42em}
\setlength{\parindent}{0pt}
\setlength{\emergencystretch}{2.5em}
\color{ink}

\titleformat{\section}
  {\sffamily\Large\bfseries\color{ink}}
  {\sffamily\Large\bfseries\color{laranja}\thesection}
  {0.7em}{}
  [\vspace{-0.55em}{\color{linhasuave}\titlerule[0.6pt]}]
\titleformat{\subsection}
  {\sffamily\large\bfseries\color{ink}}
  {\sffamily\large\bfseries\color{verdeesc}\thesubsection}
  {0.55em}{}
\titlespacing*{\section}{0pt}{1.35em}{0.65em}
\titlespacing*{\subsection}{0pt}{1.05em}{0.35em}

\pagestyle{fancy}
\fancyhf{}
\fancyhead[L]{\textcolor{muted}{\sffamily\footnotesize\nouppercase{\rightmark}}}
\fancyfoot[L]{\textcolor{muted}{\sffamily\footnotesize Computação gráfica}}
\fancyfoot[R]{\textcolor{ink}{\sffamily\small\thepage}}
\renewcommand{\headrulewidth}{0.4pt}
\renewcommand{\headrule}{%
  \hbox to\headwidth{\textcolor{linhasuave}{\leaders\hrule height \headrulewidth\hfill}}%
}
\setlength{\headheight}{24pt}
\addtolength{\topmargin}{-9.5pt}
\renewcommand{\sectionmark}[1]{\markright{\thesection\quad #1}}
\makeatletter
\patchcmd{\l@section}{1.0em \@plus\p@}{0.42em \@plus 0.2\p@}{}{\typeout{toc patch failed}}
\makeatother

\newtcblisting{codigo}[1]{
  enhanced,
  listing engine=minted,
  minted language=c,
  listing only,
  minted options={fontsize=\footnotesize, breaklines, style=default, tabsize=4},
  colback=codebg,
  colframe=codebg,
  boxrule=0pt,
  borderline west={2.4pt}{0pt}{laranja},
  sharp corners,
  left=2.6mm, right=2mm, top=1.1mm, bottom=1.1mm,
  title={#1},
  colbacktitle=codebg,
  coltitle=laranjaesc,
  fonttitle=\sffamily\scriptsize,
  attach title to upper={\par\smallskip},
  before skip=8pt,
  after skip=8pt
}

\newtcolorbox{voz}{
  enhanced, breakable,
  colback=caixafundo, colframe=caixafundo, boxrule=0pt,
  borderline west={2.4pt}{0pt}{verdeesc},
  sharp corners, left=3mm, right=2.2mm, top=1.4mm, bottom=1.4mm,
  title={Para dizer em voz alta},
  fonttitle=\sffamily\footnotesize\bfseries,
  coltitle=verdeesc, colbacktitle=caixafundo,
  attach title to upper={\par\smallskip},
  before skip=8pt, after skip=8pt
}

\newtcolorbox{conta}{
  enhanced, breakable,
  colback=contafundo, colframe=contafundo, boxrule=0pt,
  borderline west={2.4pt}{0pt}{laranja},
  sharp corners, left=3mm, right=2.2mm, top=1.4mm, bottom=1.4mm,
  title={Conta para conferir com a tecla M},
  fonttitle=\sffamily\footnotesize\bfseries,
  coltitle=laranjaesc, colbacktitle=contafundo,
  attach title to upper={\par\smallskip},
  before skip=8pt, after skip=8pt
}

\newtcolorbox{cuidado}{
  enhanced, breakable,
  colback=alertafundo, colframe=alertafundo, boxrule=0pt,
  borderline west={2.4pt}{0pt}{laranjaesc},
  sharp corners, left=3mm, right=2.2mm, top=1.4mm, bottom=1.4mm,
  title={Cuidado},
  fonttitle=\sffamily\footnotesize\bfseries,
  coltitle=laranjaesc, colbacktitle=alertafundo,
  attach title to upper={\par\smallskip},
  before skip=8pt, after skip=8pt
}

\newenvironment{figura}
  {\par\medskip\noindent\begin{minipage}{\linewidth}\centering}
  {\end{minipage}\par\medskip}

\newcommand{\fn}[1]{\texttt{\nobreak\hspace{0pt}#1}}
\newcommand{\desenhalogo}{%
\begin{tikzpicture}[scale=0.92, line cap=round, line join=round]
  \useasboundingbox (-3.45,-3.25) rectangle (3.45,3.25);
  \foreach \ang/\rot/\rx/\ry/\dist/\cor in {%
    135/45/0.64/1.16/1.42/laranjaesc,%
    45/-45/0.50/0.90/1.16/laranja,%
    225/135/0.50/0.90/1.16/verde,%
    315/225/0.64/1.16/1.42/verdeesc%
  }{%
    \begin{scope}[shift={(\ang:\dist)}, rotate=\rot]
      \foreach \f/\w in {1.00/1.05, 0.80/0.8, 0.60/0.65, 0.40/0.5, 0.22/0.4}{%
        \draw[\cor, line width=\w pt] (0,0) ellipse ({\f*\rx} and {\f*\ry});
      }
    \end{scope}
  }
\end{tikzpicture}%
}

\begin{document}
\pagenumbering{gobble}
\thispagestyle{empty}

\vspace*{0.4cm}
{\sffamily\footnotesize\color{laranjaesc} Computação gráfica \textcolor{linhasuave}{—} pipeline clássico}

\vspace{0.55cm}
{\sffamily\fontsize{26}{30}\selectfont\bfseries A logo e-colab em OpenGL\par}

\vspace{0.35cm}
{\LARGE Da janela ao tubo que desenha a pétala\par}

\vspace{0.7cm}
\noindent
\begin{minipage}[c]{0.46\textwidth}
\desenhalogo
\end{minipage}\hfill
\begin{minipage}[c]{0.50\textwidth}
\sffamily\small
Quatro cópias da mesma malha. Laranja em cima, verde embaixo. A diagonal que desce da esquerda para a direita é a maior, como na referência.

\vspace{0.7em}
O miolo de cada pétala recua em \(z\). De frente lê-se a logo. De lado, um funil.

\vspace{0.7em}
A palavra da referência não entra neste protótipo. O programa desenha o símbolo.
\end{minipage}

\vspace{1.05cm}
Esta apostila explica o programa que está na pasta \fn{examples}. A ordem é a de um quadro: primeiro a OpenGL como máquina de estados, depois a matriz que move o vértice, e só então a curva que vira tubo. Luz, tonalização e Z-buffer fecham o caminho até o pixel.

\vspace{0.45em}
Os números são os do código. Quando um comentário do fonte e a conta discordam, vale a conta, e o texto diz onde isso acontece. O caso concreto é a normal da elipse: o comentário diz ``para fora''; o vetor calculado aponta para o interior da curva. A normal que a luz enxerga é outra, e as duas coisas ficam separadas na seção da geometria.

\vspace{0.9cm}
{\color{linhasuave}\rule{\textwidth}{0.6pt}}
\vspace{0.35cm}
{\sffamily\small\color{muted}
Protótipo \fn{ecolab\_exemplo}
\quad\textcolor{muted}{·}\quad
\fn{main.c}, \fn{math3d.c}, \fn{petal.c}
\quad\textcolor{muted}{·}\quad
novembro de 2026
}

\clearpage
\pagenumbering{roman}
\pagestyle{plain}
{\setlength{\parskip}{0pt}\tableofcontents}
\clearpage
\pagenumbering{arabic}
\pagestyle{fancy}

\section{Como estudar com o programa aberto}

A leitura acompanha o executável. Cada bloco abaixo termina numa tecla. Se a tecla não produz o que a seção descreveu, a seção ainda não foi entendida: a tela é o critério.

\begin{center}
\captionof{table}{O que ler e o que apertar em seguida.}
\small
\begin{tabularx}{\linewidth}{@{}l X l@{}}
\toprule
Depois de & O que tem de estar claro & Tecla \\
\midrule
seção~\ref{sec:quadro} & quem limpa, quem projeta, quem gira & abrir o programa \\
seção~\ref{sec:petalas} & \(p' = R_z T S p\) & \fn{M}, depois \fn{+} \\
seção~\ref{sec:proj} & perspectiva muda o tamanho com a distância & \fn{P} e \fn{O} \\
seção~\ref{sec:geo} & a pétala é tubo ao longo de elipse & \fn{F}, \fn{3}, \fn{W} \\
seção~\ref{sec:luz} & Phong é a conta; Gouraud é a interpolação & \fn{L} e \fn{G} \\
seção~\ref{sec:z} & sem teste, ganha quem foi desenhado por último & \fn{Z} \\
\bottomrule
\end{tabularx}
\end{center}

No Linux, com FreeGLUT instalado:

\begin{codigo}{terminal}
sudo apt install build-essential cmake freeglut3-dev
cmake -S . -B build && cmake --build build
./build/ecolab_exemplo
\end{codigo}

A tecla \fn{M} escreve no terminal que lançou o programa, não na janela. A barra de título mostra projeção, Gouraud, luz e Z-buffer.

\begin{voz}
A cena abre inclinada de propósito. Vista frontal é a tecla \fn{F} ou \fn{1}. A referência está em \fn{docs/referencia-ecolab.jpg}; o programa não carrega essa foto. A comparação é olho no olho, com a janela ao lado.
\end{voz}

\section{A cena}
\label{sec:cena}

A referência é um símbolo plano: quatro lóbulos de traços concêntricos, laranja nos dois de cima, verde nos dois de baixo, a diagonal noroeste--sudeste um pouco maior, e a palavra e-colab por baixo. O protótipo fica com os quatro lóbulos e dá a eles volume. Cada lóbulo é a mesma malha, desenhada quatro vezes. Escala, distância ao centro e rotação mudam de uma cópia para a outra. A cor também.

\begin{figura}
\desenhalogo
\captionof{figure}{Leitura da cena. Cima, laranja; baixo, verde. Os lóbulos da diagonal principal são maiores e nascem um pouco mais longe do centro.}
\end{figura}

A malha nasce centrada na origem, comprida no eixo \(Y\). Uma translação empurra esse centro ao longo de \(+Y\). Uma rotação em \(Z\) leva esse \(+Y\) para a diagonal. O miolo da elipse não fica no plano \(z = 0\): ele recua até \(z = -0{,}50\). Por isso a tecla \fn{3}, quase de lado, mostra um funil em vez de um risco. Esse volume é o que dá trabalho à luz e ao Z-buffer. Uma elipse desenhada com \fn{GL\_LINE} no plano não teria normal nem profundidade para discutir.

O plano de estudo do grupo, em \fn{docs/PLANO.md}, ainda descreve a palavra desenhada por Bresenham ou DDA e uma placa atrás do símbolo. Esses arquivos não estão nesta árvore, e o \fn{CMakeLists.txt} não os compila. Esta apostila cobre o que o executável faz hoje. Se a palavra voltar, o raster tem de ser código do grupo: uma lista de pixels, cada pixel um quadradinho. \fn{glBegin(GL\_LINES)} na palavra não cumpre esse item.

\section{Qual OpenGL é esta}
\label{sec:ogl}

A OpenGL é uma interface de desenho. O programa descreve primitivas e estado; a implementação, dentro do driver, rasteriza. Há dois jeitos comuns de falar com ela.

No perfil moderno, a partir do OpenGL 3.3 core, o programa sobe os vértices para a memória da GPU, escreve um vertex shader que multiplica posição por matrizes, e um fragment shader que calcula a cor. \fn{glBegin}, \fn{glMatrixMode}, \fn{glLight} e \fn{glMultMatrix} não existem nesse perfil.

Este trabalho usa o outro jeito: o pipeline fixo, também chamado de modo imediato, via FreeGLUT. O vértice é emitido um a um entre \fn{glBegin} e \fn{glEnd}. A matriz entra por \fn{glMultMatrixf}. A luz é um estado da máquina, não uma função escrita por vocês. A escolha não é estética. O enunciado pede a matriz de vocês aplicada com \fn{glMultMatrix}. No perfil core essa chamada não existe.

As contas que a apostila desenvolve são as mesmas nos dois perfis. No shader moderno o vértice de recorte sai de
\[
  p_{\text{clip}} = P \, M \, p .
\]
Aqui o driver faz esse produto com a matriz de projeção e a modelview que estiverem carregadas. Mudou quem multiplica. Não mudou a álgebra.

Três famílias de função ficam de fora da entrega, mesmo que a tela ``funcione'' com elas: \fn{glTranslatef}, \fn{glRotatef}, \fn{glScalef}, \fn{glOrtho}, \fn{gluPerspective} e \fn{gluLookAt}. Elas escondem a matriz. \fn{glLoadIdentity}, \fn{glPushMatrix}, \fn{glPopMatrix} e \fn{glMatrixMode} ficam: a pilha só guarda e restaura. Quem transforma o vértice é o arranjo que vocês montaram.

\begin{cuidado}
Uma chamada proibida que acerta a imagem zera o item da matriz ou o da projeção na hora da arguição. O teste visual não distingue \fn{glRotatef} de \fn{mat4\_make\_rotation\_z}. O código distingue.
\end{cuidado}

\section{Um quadro, do clear ao swap}
\label{sec:quadro}

O FreeGLUT cria a janela e chama o programa de volta. \fn{main} registra as funções e entra em \fn{glutMainLoop}, que não retorna. Cada vez que a janela precisa ser redesenhada, o GLUT chama \fn{display}. Mudar uma variável não redesenha sozinho: \fn{glutPostRedisplay} avisa que o quadro está sujo.

\begin{codigo}{examples/main.c}
glutInitDisplayMode(GLUT_DOUBLE | GLUT_RGB | GLUT_DEPTH | GLUT_MULTISAMPLE);
glutDisplayFunc(display);
glutReshapeFunc(reshape);
glutKeyboardFunc(keyboard);
glutSpecialFunc(special_keys);
glutMouseFunc(mouse);
glutMotionFunc(motion);
glutMainLoop();
\end{codigo}

\fn{GLUT\_DOUBLE} pede dois buffers de cor. O desenho acontece no buffer de trás. \fn{glutSwapBuffers}, no fim de \fn{display}, troca os dois. Sem isso o usuário veria a limpeza e a malha pela metade. \fn{GLUT\_RGB} pede cor em vermelho, verde e azul. \fn{GLUT\_DEPTH} cria o Z-buffer. Sem esse bit, \fn{glEnable(GL\_DEPTH\_TEST)} não tem memória de profundidade e a tecla \fn{Z} não tem o que desligar. \fn{GLUT\_MULTISAMPLE} pede antialiasing nas bordas do tubo. É conforto visual, não é o algoritmo de visibilidade do trabalho.

O corpo de \fn{display} é o índice do resto da apostila.

\begin{enumerate}
\item Limpa cor e profundidade.
\item Acerta o viewport, o modo do polígono, a tonalização, o teste de profundidade e a luz, conforme as teclas.
\item \fn{glMatrixMode(GL\_PROJECTION)}, identidade, uma das duas projeções.
\item \fn{glMatrixMode(GL\_MODELVIEW)}, identidade.
\item Translação da câmera.
\item Posição da luz, ainda sem a rotação do objeto.
\item Rotação de pitch e yaw.
\item \fn{draw\_petals}, quatro vezes push, matriz, malha, pop.
\item Troca os buffers.
\end{enumerate}

\fn{glClear} não zera matriz. A identidade é chamada todo quadro, na projeção e na modelview. Sem ela, a rotação deste quadro compõe com a do quadro anterior e a cena foge.

\begin{codigo}{examples/main.c}
glClear(GL_COLOR_BUFFER_BIT | GL_DEPTH_BUFFER_BIT);
glViewport(0, 0, g_width, g_height);
glMatrixMode(GL_PROJECTION);
glLoadIdentity();
/* perspectiva ou ortogonal, e apply_matrix(proj) */
glMatrixMode(GL_MODELVIEW);
glLoadIdentity();
/* câmera, luz, pitch, yaw, draw_petals */
glutSwapBuffers();
\end{codigo}

A cor de fundo é \fn{glClearColor(0.97, 0.96, 0.94, 1)}, um branco quente. A profundidade limpa é \fn{glClearDepth(1.0)}: o valor que o buffer trata como longe. O teste, armado em \fn{init\_gl}, é \fn{GL\_LESS}. Um fragmento pinta se a profundidade dele for menor que a armazenada.

\fn{reshape} guarda a largura e a altura novas. Se a altura chegasse a zero, o aspect da projeção dividiria por zero; a função força altura mínima \(1\). O aspect é largura sobre altura. Na janela inicial, \(960/720 = 4/3\).

\section{A máquina de estados}
\label{sec:estado}

A OpenGL deste programa não recebe a câmera como argumento. Ela consulta o que estiver valendo. Cor, normal, matriz ativa, teste de profundidade, modelo de tonalização e luz são estado. Dois erros de sintoma parecido nascem daí.

O primeiro: a matriz certa foi multiplicada no modo errado. \fn{glMatrixMode(GL\_PROJECTION)} e \fn{glMatrixMode(GL\_MODELVIEW)} escolhem qual das duas matrizes \fn{glMultMatrixf} altera. Câmera aplicada com o modo ainda em projeção ``não gira nada'', porque o quadro seguinte dá \fn{glLoadIdentity} na projeção e a modelview segue identidade.

O segundo: o estado foi armado cedo demais ou tarde demais. A luz é o exemplo. \fn{glLightfv} lê a modelview corrente, transforma a posição e guarda o resultado em coordenada de olho. No quadro, isso acontece depois da translação da câmera e antes do yaw. A lâmpada fica parada enquanto a logo gira, e o brilho caminha pela superfície. A mesma chamada depois do yaw gravaria a rotação dentro da lâmpada, e o brilho grudaria na pétala.

Estes são os estados que a arguição pode apontar na tela.

\begin{center}
\captionof{table}{Estado que muda a imagem, e a chamada que escreve nele.}
\small
\begin{tabularx}{\linewidth}{@{}l X@{}}
\toprule
Estado & Onde é escrito \\
\midrule
Cor de limpeza e profundidade de limpeza & \fn{glClearColor}, \fn{glClearDepth} \\
Retângulo de pixels & \fn{glViewport} \\
Polígono preenchido ou aramado & \fn{glPolygonMode} \\
Cor por vértice ou cor por faceta & \fn{glShadeModel} \\
Teste de profundidade & \fn{glEnable}/\fn{glDisable(GL\_DEPTH\_TEST)} \\
Luz ligada & \fn{GL\_LIGHTING} \\
Qual matriz \fn{glMultMatrixf} altera & \fn{glMatrixMode} \\
Modelview e projeção & \fn{glLoadIdentity}, \fn{apply\_matrix} \\
Cópia da modelview & \fn{glPushMatrix}, \fn{glPopMatrix} \\
Cor do material & \fn{glColor3fv}, com \fn{GL\_COLOR\_MATERIAL} \\
Normal do vértice & \fn{glNormal3f} \\
\bottomrule
\end{tabularx}
\end{center}

\fn{glColor}, com a iluminação ligada, é ignorado até \fn{glEnable(GL\_COLOR\_MATERIAL)}. Sem isso a logo inteira sai cinza e o grupo procura o bug na pétala. \fn{setup\_lighting} liga o estado e manda a cor alimentar o ambiente e a difusa do material, na frente e nas costas. O especular do material não vem do \fn{glColor}; ele é fixo, declarado à parte.

\section{Do objeto à janela}
\label{sec:espacos}

Um vértice da pétala passa por espaços sucessivos. Decorá-los evita a frase ``a matriz da câmera'' para três matrizes diferentes.

\begin{enumerate}
\item \textbf{Objeto.} A elipse está centrada na origem, comprida em \(Y\). A normal do tubo também está nesse espaço. É o que \fn{petal\_draw} emite.
\item \textbf{Logo.} Escala, translação ao longo de \(Y\) e rotação em \(Z\). As quatro pétalas passam a morar em volta da origem.
\item \textbf{Mundo virado.} Yaw em \(Y\) e pitch em \(X\). A logo gira; a câmera, nesta implementação, não sai do lugar.
\item \textbf{Olho.} A translação \((0,\ 0{,}18,\ -d)\) coloca a origem da logo à frente da câmera. A câmera da OpenGL fica na origem e olha para \(-Z\).
\item \textbf{Recorte.} A projeção produz a coordenada homogênea \((x, y, z, w)\).
\item \textbf{NDC.} A divisão por \(w\) leva \(x\), \(y\) e \(z\) ao cubo \([-1, 1]\).
\item \textbf{Janela.} O viewport espalha o quadrado \([-1, 1]\times[-1, 1]\) nos pixels. No mundo da OpenGL, \(y\) cresce para cima. No mouse do GLUT, \(y\) cresce para baixo.
\item \textbf{Fragmento.} Cada pixel coberto ganha cor e profundidade. O Z-buffer decide se essa profundidade substitui a que já estava lá.
\end{enumerate}

\begin{figura}
\begin{tikzpicture}[
  caixa/.style={draw=linhasuave, fill=codebg, rounded corners=1.5pt,
                font=\sffamily\footnotesize, align=center, inner sep=4.5pt,
                minimum height=8mm, text=ink},
  seta/.style={-{Stealth[length=2mm]}, semithick, draw=muted}
]
\node[caixa] (a) {objeto};
\node[caixa, right=3.2mm of a] (b) {logo\\[-0.15em]{\scriptsize\color{muted} \(R_z T S\)}};
\node[caixa, right=3.2mm of b] (c) {olho\\[-0.15em]{\scriptsize\color{muted} \(T_{\mathrm{cam}} R_x R_y\)}};
\node[caixa, right=3.2mm of c] (d) {recorte\\[-0.15em]{\scriptsize\color{muted} \(P\)}};
\node[caixa, right=3.2mm of d] (e) {pixel};
\draw[seta] (a) -- (b);
\draw[seta] (b) -- (c);
\draw[seta] (c) -- (d);
\draw[seta] (d) -- (e);
\end{tikzpicture}
\captionof{figure}{Cinco nomes para não chamar tudo de ``a matriz''. A projeção \(P\) não é multiplicada por vocês com a modelview: o driver faz \(P\,M\,p\).}
\end{figura}

Há dois sistemas destros, e os dois são verdade ao mesmo tempo. No mundo, antes da câmera, \(+X\) aponta para a direita da tela, \(+Y\) para cima e \(+Z\) para quem olha. A câmera da OpenGL, porém, senta na origem e olha para o eixo \(Z\) negativo. Um ponto à frente da câmera tem \(z\) de olho negativo. A translação da câmera existe para fabricar esse \(z\) negativo: ela empurra a logo para \(z = -d\).

\begin{voz}
Ponto à frente da câmera tem \(z\) negativo no espaço do olho. O fundo do Z-buffer é \(1\), que corresponde ao lado longe do NDC. Perto é profundidade pequena. \fn{GL\_LESS} guarda o menor valor, isto é, quem está mais perto.
\end{voz}

\section{Coordenada homogênea}
\label{sec:homog}

Translação não é transformação linear. \(T(p) = p + t\) não passa pela origem, então não cabe numa matriz \(3\times 3\). A coordenada homogênea acrescenta uma quarta componente e devolve a translação ao mundo das matrizes.

Um ponto do espaço vira a coluna \((x, y, z, 1)\). Um vetor, que não tem posição, vira \((x, y, z, 0)\). A mesma matriz \(4\times 4\) age nos dois. Na translação, a última coluna é \((t_x, t_y, t_z, 1)\). Multiplicando um ponto, o \(1\) puxa \(t\) para o resultado. Multiplicando um vetor, o \(0\) anula \(t\): direção não se translada. Normal é direção. Por isso a translação da pétala não empurra a normal, e ainda assim a rotação e a escala atuam nela.

A quarta linha da matriz de projeção perspectiva não é \((0, 0, 0, 1)\). Ela escreve \(-z\) em \(w\). Depois da matriz, a GPU divide \(x\), \(y\) e \(z\) por \(w\). Essa divisão é a perspectiva. \fn{mat4\_transform\_point} não a faz. A tecla \fn{M} mostra um ponto no espaço da logo, ainda sem projeção e sem câmera.

A luz usa a mesma distinção. A posição em \fn{display} é \((1{,}8,\ 2{,}4,\ 5{,}0,\ 1)\). O \(1\) final declara luz pontual: a OpenGL transforma esse valor como ponto. Se a quarta componente fosse \(0\), seria luz direcional, transformada como vetor, e os raios chegariam paralelos.

\section{Como a matriz mora na memória}
\label{sec:memoria}

No papel a matriz se escreve por linhas. Na OpenGL ela se armazena por colunas. O elemento da linha \(r\) e da coluna \(c\), com índices começando em zero, fica em
\begin{equation}
  m[c \cdot 4 + r].
  \label{eq:indice}
\end{equation}
A translação, que no papel está na última coluna, ocupa os índices \(12\), \(13\) e \(14\). Não ocupa o fim de cada linha.

\begin{figura}
\begin{tikzpicture}[
  cel/.style={draw=linhasuave, minimum width=1.48cm, minimum height=0.76cm,
              outer sep=0pt, inner sep=0pt, font=\small, anchor=center},
  rot/.style={font=\sffamily\scriptsize, text=muted, align=center}
]
\foreach \c in {0,1,2,3} {
  \foreach \r in {0,1,2,3} {
    \pgfmathtruncatemacro{\idx}{\c*4+\r}
    \ifnum\c=3
      \node[cel, fill=contafundo] at (\c*1.48, -\r*0.76) {\idx};
    \else
      \node[cel, fill=white] at (\c*1.48, -\r*0.76) {\idx};
    \fi
  }
}
\node[rot] at (0.00, 0.78) {coluna 0};
\node[rot] at (1.48, 0.78) {coluna 1};
\node[rot] at (2.96, 0.78) {coluna 2};
\node[rot] at (4.44, 0.78) {translação};
\draw[-{Stealth[length=1.8mm]}, laranjaesc, semithick]
  (-1.15, 0.28) -- (-1.15, -2.55);
\node[rot, rotate=90] at (-1.72, -1.14) {a memória desce a coluna};
\end{tikzpicture}
\captionof{figure}{Os dezesseis endereços. A coluna da translação está destacada. \(t_y\), por exemplo, é \(m[13]\).}
\end{figura}

\fn{mat4\_identity} zera o arranjo e escreve \(1\) em \(m[0]\), \(m[5]\), \(m[10]\) e \(m[15]\), a diagonal. \fn{mat4\_print} percorre linha por linha na hora de imprimir, lendo \(m[c\cdot 4+r]\), para o terminal mostrar a matriz do jeito do caderno. A memória por baixo continua por coluna.

A multiplicação segue a mesma convenção. \fn{mat4\_multiply(out, a, b)} calcula \(a\,b\), com a coluna à direita.

\begin{codigo}{examples/math3d.c · o elemento (linha, coluna)}
r[col * 4 + row] =
    a[0 * 4 + row] * b[col * 4 + 0] +
    a[1 * 4 + row] * b[col * 4 + 1] +
    a[2 * 4 + row] * b[col * 4 + 2] +
    a[3 * 4 + row] * b[col * 4 + 3];
\end{codigo}

O índice \fn{col * 4 + row} é a equação~\eqref{eq:indice}. O somatório é a regra de sempre: a linha de \(a\) contra a coluna de \(b\). O resultado vai primeiro para um arranjo temporário, e só depois para \fn{out}. Por isso \fn{out} pode ser o próprio \fn{a} ou o próprio \fn{b} sem que a conta se pise.

\begin{codigo}{examples/math3d.c · ponto coluna, com w = 1}
*ox = m[0] * x + m[4] * y + m[8] * z + m[12];
*oy = m[1] * x + m[5] * y + m[9] * z + m[13];
*oz = m[2] * x + m[6] * y + m[10] * z + m[14];
\end{codigo}

Cada linha lê uma linha da matriz, descendo quatro colunas. É \(M\) vezes \((x, y, z, 1)\). A componente \(w\) do resultado não é calculada: para as matrizes afins desta cena ela continua \(1\).

\begin{cuidado}
Montar a matriz por linha e entregar esse vetor a \fn{glMultMatrixf} coloca \(t_x\) em \(m[3]\), que a OpenGL lê como o fim da primeira coluna, não como translação. A imagem entorta de um jeito que não se conserta chutando o ângulo.
\end{cuidado}

\section{glMultMatrixf multiplica pela direita}
\label{sec:ordem}

\fn{apply\_matrix} é a única porta das matrizes do grupo.

\begin{codigo}{examples/math3d.c}
void apply_matrix(const float m[16])
{
    glMultMatrixf(m);
}
\end{codigo}

\fn{glMultMatrixf(A)} substitui a matriz corrente \(M\) por \(M A\). O vértice é coluna à direita, então a última matriz aplicada é a primeira que toca o vértice. Depois de \fn{glLoadIdentity}:
\begin{equation}
  \texttt{apply}(A);\; \texttt{apply}(B);\; \texttt{apply}(C);
  \quad\Longleftrightarrow\quad
  p' = A\, B\, C\, p .
  \label{eq:ordem}
\end{equation}
\(C\) age primeiro, \(B\) em seguida, \(A\) por último. A frase para levar à arguição é esta: a chamada de baixo, no fonte, é a transformação de dentro, na álgebra.

A pilha não entra nessa conta. \fn{glPushMatrix} copia a modelview. \fn{glPopMatrix} restaura a cópia. Nenhuma das duas gira vértice. Sem o pop, a segunda pétala herda \(R_z T S\) da primeira e a rotação se acumula.

\section{Translação, escala e rotação}
\label{sec:afim}

As seis funções de \fn{math3d.c} preenchem uma identidade e escrevem o bloco que muda. O ângulo entra em radianos. Em \fn{main.c} os graus são convertidos na hora da chamada, multiplicando por \(\pi/180\). Dentro de \fn{petal.c} o parâmetro da elipse já nasce em radianos: uma volta é \(2\pi\).

\subsection{Translação e escala}

\[
T(t) =
\begin{pmatrix}
1 & 0 & 0 & t_x \\
0 & 1 & 0 & t_y \\
0 & 0 & 1 & t_z \\
0 & 0 & 0 & 1
\end{pmatrix}
\qquad
S(s) =
\begin{pmatrix}
s_x & 0 & 0 & 0 \\
0 & s_y & 0 & 0 \\
0 & 0 & s_z & 0 \\
0 & 0 & 0 & 1
\end{pmatrix}
\]

No arranjo, \(t_x, t_y, t_z\) vão para \(m[12]\), \(m[13]\), \(m[14]\). A escala vai para \(m[0]\), \(m[5]\), \(m[10]\). Escala uniforme, o caso da pétala, usa o mesmo número nos três. Escala com \(s_x \neq s_y\) mudaria a forma e também a direção da normal; a deste programa não faz isso.

\subsection{As três rotações}

A rotação positiva segue a mão direita: o polegar aponta no sentido do eixo, e os dedos fecham no sentido do ângulo positivo. Na tela, a rotação positiva em \(Z\) é anti-horária. É a que posiciona as pétalas.

\[
R_z =
\begin{pmatrix}
c & -s & 0 & 0 \\
s &  c & 0 & 0 \\
0 &  0 & 1 & 0 \\
0 &  0 & 0 & 1
\end{pmatrix}
\quad
R_x =
\begin{pmatrix}
1 & 0 & 0 & 0 \\
0 & c & -s & 0 \\
0 & s &  c & 0 \\
0 & 0 & 0 & 1
\end{pmatrix}
\quad
R_y =
\begin{pmatrix}
c & 0 & s & 0 \\
0 & 1 & 0 & 0 \\
-s & 0 & c & 0 \\
0 & 0 & 0 & 1
\end{pmatrix}
\]

com \(c = \cos\theta\) e \(s = \sin\theta\). A conferência com o fonte é pelo índice. Em \fn{mat4\_make\_rotation\_z}: \(m[0] = c\), \(m[1] = s\), \(m[4] = -s\), \(m[5] = c\). Quem escrever \(m[1] = -s\) espelha a logo.

\(R_x\) gira \(Y\) na direção de \(Z\). \(R_y\) gira \(Z\) na direção de \(X\), e portanto leva um ponto em \(+X\) para um \(z\) negativo. As duas aparecem na câmera, não na pétala.

\subsection{O que a escala faz com a normal}

A OpenGL transforma a normal pela inversa transposta do bloco \(3\times 3\) da modelview. A translação sai da conta, porque normal é vetor. A razão da inversa transposta é geométrica: a normal tem de continuar perpendicular à tangente depois da transformação.

Se uma tangente \(t\) vira \(A t\), a normal \(n'\) precisa cumprir \(n' \cdot (A t) = n \cdot t\) para toda tangente. Isso força \(n' = (A^{-1})^{\mathsf T} n\).

Numa escala não uniforme a direção muda. Um caso pequeno basta. A tangente \((1, 1)\) tem normal \((-1, 1)\). Depois de \(A = \mathrm{diag}(2, 1)\), a tangente vira \((2, 1)\). A normal transformada pela inversa transposta é \((-1/2,\ 1)\), e o produto escalar com \((2, 1)\) continua zero. A normal ingênua \(A n = (-2, 1)\) não é mais perpendicular.

Na pétala a escala é uniforme, \(S = s I\). A inversa transposta vale \((1/s) I\): a direção sobrevive e o comprimento vira \(1/s\). A malha já entrega normal de comprimento \(1\), feita por \fn{normalize3}. A modelview desfaz esse comprimento sempre que \(s \neq 1\). A pétala grande usa \(s = 1{,}10\), a pequena usa \(s = 0{,}84\), e a tecla \fn{+} muda outro fator por cima. Sem uma renormalização, \(N\cdot L\) sairia em escala errada: a pétala grande ficaria escura demais, a pequena clara demais, e o brilho especular mudaria de largura.

\fn{glEnable(GL\_NORMALIZE)} faz essa segunda normalização, já no espaço do olho, para toda normal que a modelview tocar. As duas normalizações têm empregos diferentes. \fn{normalize3} acerta o tubo no espaço do objeto. \fn{GL\_NORMALIZE} desfaz o estrago da escala.

\section{As quatro pétalas}
\label{sec:petalas}

\begin{codigo}{examples/main.c · uma cópia}
glPushMatrix();
mat4_make_rotation_z(rz, degrees[i] * MATH3D_PI / 180.f);
apply_matrix(rz);
mat4_make_translation(t, 0.f, out, 0.f);
apply_matrix(t);
mat4_make_scale(s, size, size, size);
apply_matrix(s);
petal_draw(/* laranja se i < 2, senão verde */);
glPopMatrix();
\end{codigo}

Pela equação~\eqref{eq:ordem},
\begin{equation}
  p' = R_z(\theta)\, T(0, d, 0)\, S(s)\, p .
  \label{eq:petala}
\end{equation}
A escala acontece em volta da origem local da pétala. A translação empurra essa origem ao longo de \(+Y\), por uma distância que a escala não altera. A rotação leva o eixo \(+Y\) para a diagonal. Essa separação é o conteúdo da tecla \fn{+}: os anéis crescem, o centro da pétala não se afasta. Com escala alta, o lóbulo avança sobre a origem e as quatro malhas se atravessam. É o comportamento certo desta ordem, não um acidente.

Os quatro ângulos e os dois tamanhos:

\begin{center}
\captionof{table}{As quatro chamadas de \fn{petal\_draw}. Ângulo positivo é anti-horário. ``Grande'' usa escala \(1{,}10\) e distância \(0{,}44\times 1{,}04\). ``Pequena'' usa \(0{,}84\) e \(0{,}44\times 0{,}90\).}
\small
\begin{tabular}{@{}r r l l l@{}}
\toprule
\(i\) & \(\theta\) & lugar & cor & tamanho \\
\midrule
0 & \(-45^\circ\) & superior direito & laranja & pequeno \\
1 & \(+45^\circ\) & superior esquerdo & laranja & grande \\
2 & \(135^\circ\) & inferior esquerdo & verde & pequeno \\
3 & \(225^\circ\) & inferior direito & verde & grande \\
\bottomrule
\end{tabular}
\end{center}

A conta do lugar: a translação produz o ponto \((0, d, 0)\), e \(R_z\) leva esse ponto a \((-d\sin\theta,\ d\cos\theta,\ 0)\). Para \(\theta = +45^\circ\), seno e cosseno valem \(\sqrt{2}/2\), então o centro cai em \(x\) negativo e \(y\) positivo: superior esquerdo. Para \(\theta = -45^\circ\), o seno troca de sinal e o centro vai para o superior direito. A diagonal grande é a dos índices \(1\) e \(3\), noroeste e sudeste, que é a diagonal mais cheia da referência.

A elipse é comprida em \(Y\) e está centrada na própria origem. Metade dela aponta de volta para o centro da logo. No lóbulo grande, o centro fica a \(0{,}458\) da origem e a semi-altura, já escalada, vale cerca de \(0{,}605\). A ponta interna cai em \(0{,}458 - 0{,}605 = -0{,}147\): atravessa a origem. A ponta externa vai a \(1{,}06\). No lóbulo pequeno a ponta interna fica em \(-0{,}066\). O miolo da logo é zona de cruzamento de propósito. Quem pinta por cima, nessa zona, é decisão do Z-buffer.

\begin{conta}
Escala da tecla \fn{M} igual a \(1\), translação \((0,\ 0{,}44,\ 0)\), rotação de \(45^\circ\), ponto local \((0{,}20,\ 0,\ 0)\).

\(S p = (0{,}20,\ 0,\ 0)\).

\(T S p = (0{,}20,\ 0{,}44,\ 0)\).

\(c = s = 0{,}707107\).

\begin{align*}
x' &= c\cdot 0{,}20 - s\cdot 0{,}44 = -0{,}1697,\\
y' &= s\cdot 0{,}20 + c\cdot 0{,}44 = 0{,}4525.
\end{align*}

O \fn{printf} com três casas mostra \fn{(-0.170, 0.453, 0.000)}. A matriz \(R_z T S\), impressa por linha, tem a translação já rotacionada na última coluna:
\[
\begin{pmatrix}
0{,}707 & -0{,}707 & 0 & -0{,}311 \\
0{,}707 &  0{,}707 & 0 &  0{,}311 \\
0 & 0 & 1 & 0 \\
0 & 0 & 0 & 1
\end{pmatrix}
\]
Com a escala em \(1{,}2\) o mesmo ponto cai em \((-0{,}141,\ 0{,}481,\ 0)\). A coluna da translação não muda: \(S\) entrou antes de \(T\).
\end{conta}

\begin{cuidado}
A tecla \fn{M} não imprime a matriz de uma pétala da cena. Ela usa a escala global sozinha, a distância \(0{,}44\) sem o fator \(1{,}04\) ou \(0{,}90\), e o ângulo fixo de \(45^\circ\). É a conta limpa do caderno. A pétala de \(+45^\circ\) que está na tela ainda multiplica por \(1{,}10\) e afasta o centro até \(0{,}458\).
\end{cuidado}

Vale a pena prever a ordem trocada antes de experimentar. Se as chamadas passarem a ser \fn{apply(S)}, \fn{apply(T)}, \fn{apply(Rz)}, o produto vira \(S\, T\, R_z\). A rotação age com a pétala ainda centrada, e girar em volta da origem não tira o centro do lugar. A translação seguinte manda as quatro cópias para o mesmo ponto \((0, d, 0)\). A escala, agora por fora, cresce em volta da origem do mundo e aumenta também a distância. As quatro pétalas empilham no eixo \(Y\) e a tecla \fn{+} passa a afastá-las. A ordem do fonte não é convenção de escrita.

\begin{figura}
\begin{tikzpicture}[>=Stealth, line cap=round]
  \newcommand{\eixosordem}{%
    \draw[linhasuave] (-0.62,0) -- (0.46,0);
    \draw[linhasuave] (0,-0.48) -- (0,0.86);
    \draw[->, muted] (0,0) -- (0.40,0);
    \draw[->, muted] (0,0) -- (0,0.78);
  }
  \begin{scope}[x=3.55cm, y=3.55cm]
    \eixosordem
    \draw[muted] (0,0) ellipse (0.20 and 0.34);
    \fill (0.20,0) circle (0.016);
    \node[font=\sffamily\scriptsize, anchor=north] at (0,-0.56) {1\quad local};
  \end{scope}
  \begin{scope}[xshift=5.05cm, x=3.55cm, y=3.55cm]
    \eixosordem
    \draw[laranja!80] (0,0.44) ellipse (0.20 and 0.34);
    \fill[laranja!80] (0.20,0.44) circle (0.016);
    \node[font=\sffamily\scriptsize, anchor=north] at (0,-0.56) {2\quad depois de \(T\)};
  \end{scope}
  \begin{scope}[xshift=10.1cm, x=3.55cm, y=3.55cm]
    \eixosordem
    \begin{scope}[shift={(-0.311,0.311)}, rotate=45]
      \draw[laranjaesc, thick] (0,0) ellipse (0.20 and 0.34);
    \end{scope}
    \fill[laranjaesc] (-0.170,0.453) circle (0.018);
    \node[font=\sffamily\scriptsize, anchor=north] at (0,-0.56) {3\quad depois de \(R_z(45^\circ)\)};
  \end{scope}
\end{tikzpicture}
\captionof{figure}{A mesma elipse, com escala \(1\), nos três tempos da equação~\eqref{eq:petala}. Os três quadros usam a mesma escala. O ponto é o \((0{,}20,\ 0)\) da tecla \fn{M}: primeiro na ponta direita, depois em \((0{,}20,\ 0{,}44)\), por fim em \((-0{,}17,\ 0{,}45)\).}
\end{figura}

\section{A câmera}
\label{sec:camera}

Não há \fn{gluLookAt}. A câmera é uma translação e duas rotações, as mesmas funções da pétala.

\begin{codigo}{examples/main.c}
mat4_make_translation(cam, 0.f, 0.18f, -g_cam_dist);
apply_matrix(cam);
glLightfv(GL_LIGHT0, GL_POSITION, light_pos);
mat4_make_rotation_x(rx, g_pitch * MATH3D_PI / 180.f);
apply_matrix(rx);
mat4_make_rotation_y(ry, g_yaw * MATH3D_PI / 180.f);
apply_matrix(ry);
\end{codigo}

O vértice da logo, já transformado pela equação~\eqref{eq:petala}, chega ao olho por
\begin{equation}
  p_{\text{olho}} = T(0,\ 0{,}18,\ -d)\, R_x(\text{pitch})\, R_y(\text{yaw})\, p_{\text{logo}} .
  \label{eq:camera}
\end{equation}
Primeiro o yaw, em torno de \(Y\). Depois o pitch, em torno de \(X\). Por último o recuo. A câmera permanece na origem do olho, olhando para \(-Z\). O que se move é a logo. Para um objeto na origem, girar o objeto diante de uma câmera fixa tem o mesmo enquadramento que orbitar a câmera. A luz desfaz a ilusão: ela é gravada antes das rotações, então não gira junto. O par câmera--luz fica rígido, e a superfície da logo passa por baixo do brilho. É um prato giratório com lâmpada fixa, não uma câmera andando em volta de uma lâmpada pregada na parede.

O \(0{,}18\) em \(Y\) coloca a origem da logo um pouco acima do centro da janela. A distância \(d\) começa em \(5{,}2\), e a roda do mouse a varia entre \(1{,}6\) e \(16\). Aumentar \(d\) afasta. Na perspectiva a logo encolhe porque foi para mais longe. Na projeção ortográfica ela encolhe por outro motivo, tratado na seção seguinte: o retângulo do mundo cresce com \(d\).

O yaw inicial é \(-32^\circ\) e o pitch inicial é \(26^\circ\). A vista de cima, tecla \fn{4}, usa pitch \(80^\circ\). Isso confirma o sinal: pitch positivo traz o \(+Y\) do modelo para o lado \(+Z\), isto é, para perto da câmera, e a janela passa a mostrar a parte de cima. O pitch fica preso a \(\pm 89^\circ\). Em \(\pm 90^\circ\) o eixo do yaw alinha com a direção de visada e a vista vira de cabeça para baixo ao cruzar o polo.

Juntando as duas fórmulas, um vértice local da malha percorre
\begin{equation}
  p_{\text{olho}}
  = T_{\text{cam}}\, R_x\, R_y\, R_z\, T\, S\, p .
  \label{eq:cadeia}
\end{equation}
A projeção não entra nesse produto. O driver, no pipeline fixo, completa
\[
  p_{\text{clip}} = P\, p_{\text{olho}} .
\]

\begin{voz}
A luz é declarada depois de \(T_{\text{cam}}\) e antes de \(R_x\) e \(R_y\). Espaço ligado, a logo gira e o brilho anda. Se \fn{glLightfv} escorregar para baixo das rotações, o brilho gruda na pétala e a demonstração de Phong perde o movimento.
\end{voz}

\section{As duas projeções}
\label{sec:proj}

As duas matrizes são montadas em \fn{math3d.c} e aplicadas depois de \fn{glMatrixMode(GL\_PROJECTION)} e \fn{glLoadIdentity}. A tecla \fn{P} escolhe a perspectiva. A tecla \fn{O} escolhe a ortográfica. O código chama a segunda de ortogonal; o nome de livro é projeção ortográfica. São a mesma transformação: raios paralelos, sem ponto de fuga.

\subsection{Perspectiva}

O campo de visão vertical é \(\varphi = 40^\circ\). Seja \(a\) o aspect, largura sobre altura, e sejam \(n\) e \(f\) os planos perto e longe. No programa, \(n = 0{,}1\) e \(f = 80\).
\begin{equation}
  f_{\!v} = \frac{1}{\tan(\varphi/2)},
  \label{eq:fovy}
\end{equation}
\begin{equation}
P =
\begin{pmatrix}
f_{\!v}/a & 0 & 0 & 0 \\
0 & f_{\!v} & 0 & 0 \\
0 & 0 & A & B \\
0 & 0 & -1 & 0
\end{pmatrix}
\quad
\begin{aligned}
A &= \dfrac{f+n}{n-f},\\[0.35em]
B &= \dfrac{2fn}{n-f}.
\end{aligned}
\label{eq:persp}
\end{equation}

No arranjo: \(m[0] = f_{\!v}/a\), \(m[5] = f_{\!v}\), \(m[10] = A\), \(m[11] = -1\), \(m[14] = B\). O \(-1\) está na última linha, terceira coluna. Ele copia \(-z_{\text{olho}}\) para \(w\).

Um ponto de olho \((x, y, z, 1)\), com \(z < 0\), sai com \(w' = -z\). A divisão perspectiva produz
\[
  x_{\mathrm{ndc}} = \frac{f_{\!v}\, x}{a\,(-z)},
  \qquad
  y_{\mathrm{ndc}} = \frac{f_{\!v}\, y}{-z}.
\]
Quanto mais longe, maior \(-z\), menor o ponto na tela. Essa é a frase inteira da perspectiva. O tamanho na janela é uma razão entre o tamanho no mundo e a distância.

Na distância padrão \(d = 5{,}2\), \(\tan 20^\circ \approx 0{,}364\) e a meia-altura visível no plano da origem vale \(5{,}2 \times 0{,}364 \approx 1{,}89\). A ponta externa da pétala grande fica perto de \(1{,}06\) do centro, e na diagonal isso projeta cerca de \(0{,}75\) em \(Y\). Cabe, com folga, no campo de \(40^\circ\).

Os valores numéricos de \(A\) e \(B\), com \(n = 0{,}1\) e \(f = 80\), são \(A \approx -1{,}003\) e \(B \approx -0{,}200\). A linha de \(z\) não é uma escala uniforme. Depois da divisão por \(w\), passos iguais de \(z\) no mundo viram passos desiguais no buffer. Perto da câmera a precisão é maior. Com \(f/n = 800\) isso já seria um problema numa cena enorme; aqui o objeto inteiro mora numa faixa estreita de \(z\) e não há o que brigar com o buffer.

Trocar \(\varphi\) para \(20^\circ\), sem mexer em \(d\), troca \(\tan 20^\circ\) por \(\tan 10^\circ \approx 0{,}176\). A meia-altura visível cai para cerca de \(0{,}92\). A logo ocupa bem mais da janela e continua dentro do quadro. É o exercício clássico: campo menor, objeto maior, e a câmera só precisa recuar se a ponta passar de \(\pm 1\) em NDC.

\subsection{Ortográfica}

\begin{equation}
O =
\begin{pmatrix}
\dfrac{2}{r-l} & 0 & 0 & -\dfrac{r+l}{r-l} \\[0.7em]
0 & \dfrac{2}{t-b} & 0 & -\dfrac{t+b}{t-b} \\[0.7em]
0 & 0 & \dfrac{-2}{f-n} & -\dfrac{f+n}{f-n} \\[0.7em]
0 & 0 & 0 & 1
\end{pmatrix}
\label{eq:orto}
\end{equation}

A última linha é \((0, 0, 0, 1)\). Então \(w\) continua \(1\), a divisão perspectiva não altera tamanho, e um objeto longe sai com o mesmo tamanho de um objeto perto. Retas paralelas continuam paralelas. Não há ponto de fuga.

A chamada usa um retângulo centrado, de meia-altura \(h = 0{,}45\, d\), meia-largura \(h \cdot a\), e a faixa de profundidade de \(-20\) a \(20\):
\[
  \fn{mat4\_make\_ortho}(-h a,\ h a,\ -h,\ h,\ -20,\ 20).
\]
Com \(d = 5{,}2\), \(h = 2{,}34\). O \(x\) e o \(y\) do centro somem, porque \(r = -l\) e \(t = -b\): a translação da projeção é zero e a origem do olho cai no centro da janela. A faixa \([-20, 20]\) é larga de propósito. O parâmetro perto negativo é legal nesta fórmula; ele só diz que um dos planos fica atrás do olho. A logo, com \(d\) entre \(1{,}6\) e \(16\), permanece dentro da fatia, então o Z-buffer ainda recebe uma ordem de profundidade.

Na mesma pose, as duas projeções não enquadram igual. A perspectiva, a \(5{,}2\) de distância, mostra cerca de \(\pm 1{,}89\) em \(Y\). A ortográfica mostra \(\pm 2{,}34\). A logo aparece um pouco menor com a tecla \fn{O}. O fator \(0{,}45\) não foi escolhido para igualar os dois quadros. Ele faz a roda do mouse valer também na ortográfica: aumentar \(d\) aumenta o retângulo, e a logo encolhe. O gesto de zoom é o mesmo. O mecanismo, não.

\begin{voz}
Perspectiva: \(w = -z\), e o longe encolhe. Ortográfica: \(w = 1\), e o longe conserva o tamanho. As duas fotos da arguição saem da mesma pose. Quem mexe na seta entre um \fn{P} e um \fn{O} estraga a comparação.
\end{voz}

No Windows os nomes \fn{near} e \fn{far} são macros vazias em \fn{windows.h}. Os parâmetros do fonte se chamam \fn{znear} e \fn{zfar}. \fn{gl\_platform.h} existe para incluir o cabeçalho do Windows sem deixar essas macros comerem a matemática.

\section{A geometria da pétala}
\label{sec:geo}

A pétala não é uma superfície fechada com miolo preenchido. É um feixe de traços, e cada traço tem volume: um tubo fino ao longo de uma elipse. Volume dá normal. Normal dá o que a luz calcular.

Quatro inteiros, no topo de \fn{petal.c}, mandam na malha.

\begin{center}
\captionof{table}{Resolução da malha. Mexer num deles muda uma coisa só; a tabela da seção~\ref{sec:constantes} diz qual.}
\small
\begin{tabular}{@{}l r p{8.6cm}@{}}
\toprule
Nome & Valor & Papel \\
\midrule
\fn{PETAL\_RINGS} & 20 & traços do miolo até a borda, em cada mecha \\
\fn{PETAL\_STRANDS} & 3 & famílias de elipses, com expoentes diferentes \\
\fn{PETAL\_SEGMENTS} & 96 & amostras de uma volta \\
\fn{PETAL\_SLICES} & 8 & lados da seção do tubo \\
\bottomrule
\end{tabular}
\end{center}

São \(20 \times 3 = 60\) tubos por pétala, quatro pétalas, cada tubo com oito fitas. O modo imediato emite isso de novo a cada quadro. Para uma logo que cabe na tela, o custo é irrelevante. O que importa é ler a malha.

\subsection{A elipse que também anda em z}

\fn{ellipse\_frame} devolve um ponto da curva e dois eixos da seção.

\begin{codigo}{examples/petal.c · ponto e tangente}
*px = rx * c;
*py = ry * s;
*pz = z0 + (z1 - z0) * (t / (2.f * MATH3D_PI));
float tx = -rx * s;
float ty = ry * c;
float tz = (z1 - z0) / (2.f * MATH3D_PI);
\end{codigo}

Em fórmula, com \(t\) de \(0\) a \(2\pi\),
\begin{equation}
\begin{aligned}
p(t) &= \big(r_x\cos t,\ r_y\sin t,\ z_0 + (z_1 - z_0)\, t/(2\pi)\big),\\[0.25em]
T_{\text{cru}} &= \big(-r_x\sin t,\ r_y\cos t,\ (z_1 - z_0)/(2\pi)\big).
\end{aligned}
\label{eq:elipse}
\end{equation}
O ponto \((p_x, p_y)\) é a elipse de sempre. O \(z\) sobe linearmente com \(t\): numa volta, a curva vai de \(z_0\) a \(z_1\). De frente, o observador olha contra \(Z\) e vê uma elipse. De lado, vê a rampa. A tangente é a derivada, e em seguida \fn{normalize3} a põe com comprimento \(1\). Se o vetor fosse nulo, a função devolveria \((0, 0, 1)\) para não produzir um \(NaN\) lá na frente, na divisão pelo comprimento.

\(t = 0\) é o ponto direito da elipse, \((r_x, 0)\). A tangente no plano \(XY\), antes do \(z\), aponta para \(+Y\). O parâmetro cresce no sentido anti-horário.

\subsection{O triedro da seção}

O tubo precisa de um plano perpendicular à tangente, e de dois eixos nesse plano, para apoiar o círculo da seção.

\begin{codigo}{examples/petal.c · normal de seção e binormal}
*nx = -ty;
*ny = tx;
*nz = 0.f;
normalize3(nx, ny, nz);
*bx = ty * (*nz) - tz * (*ny);
*by = tz * (*nx) - tx * (*nz);
*bz = tx * (*ny) - ty * (*nx);
normalize3(bx, by, bz);
\end{codigo}

A tangente já está normalizada quando essas linhas rodam. O vetor \(N\) nasce de girar a parte \(XY\) da tangente \(90^\circ\) no sentido anti-horário, com componente \(z\) nula:
\begin{equation}
  N_{\text{cru}} = (-\hat T_y,\ \hat T_x,\ 0),
  \qquad
  \hat N = N_{\text{cru}} / \|N_{\text{cru}}\|.
  \label{eq:normal-secao}
\end{equation}
No ponto direito, \(\hat T \approx (0, 1, 0)\) e \(\hat N \approx (-1, 0, 0)\). A normal aponta para a origem da elipse, para o interior, não para fora. O comentário do fonte diz o contrário. A conta é a equação~\eqref{eq:normal-secao}. Para uma curva percorrida no sentido anti-horário, girar a tangente \(90^\circ\) anti-horário aponta para a esquerda de quem caminha na curva, e a esquerda é o lado de dentro.

Isso não estraga o tubo. \(N\) aqui não é a normal de iluminação. É um eixo da seção. A binormal completa um triedro destro:
\begin{equation}
  \hat B = \hat T \times \hat N .
  \label{eq:binormal}
\end{equation}
O produto do fonte é exatamente esse, componente a componente. Como \(N\) foi construído perpendicular a \(T\) --- o produto \(\hat T_x (-\hat T_y) + \hat T_y \hat T_x\) se anula, e \(N_z = 0\) --- o par \(N, B\) forma uma base do plano normal à curva. No ponto direito, com a rampa de \(z\) ainda pequena, \(B\) aponta para \(+Z\), para o lado da câmera.

Quem já viu o triedro de Frenet no cálculo reconhece as três letras e o produto vetorial. A normal daqui não é a normal principal de curvatura da curva no espaço. É a normal da elipse no plano \(XY\), reaproveitada como eixo. Com a rampa de \(z\) pequena, as duas ficam próximas. Quando a rampa é íngreme, deixam de coincidir, e o tubo continua bem definido porque \(B = T \times N\) mantém a seção perpendicular a \(T\).

\subsection{O tubo}

\begin{codigo}{examples/petal.c}
float c = cosf(angle);
float s = sinf(angle);
float ox = c * nx + s * bx;
float oy = c * ny + s * by;
float oz = c * nz + s * bz;
glNormal3f(ox, oy, oz);
glVertex3f(px + radius * ox, py + radius * oy, pz + radius * oz);
\end{codigo}

\begin{equation}
  o(\alpha) = \cos\alpha\;\hat N + \sin\alpha\;\hat B,
  \qquad
  v = p + \rho\, o(\alpha),
  \qquad
  n = o(\alpha).
  \label{eq:tubo}
\end{equation}
\(\alpha\) dá a volta na seção. \(\rho\) é \fn{PETAL\_TUBE\_RADIUS}, \(0{,}0082\). O vértice está a \(\rho\) do eixo da curva, na direção \(o\). A normal enviada à OpenGL é a mesma direção. Ela aponta para fora do fio, em todos os lados: para o interior da elipse quando \(\alpha = 0\), para o exterior quando \(\alpha = \pi\), e para os lados conforme \(B\). É essa normal que entra em \(N\cdot L\).

\begin{figura}
\begin{tikzpicture}[scale=1.55, >=Stealth]
  \draw[->, muted] (-1.85,0) -- (1.95,0) node[right, font=\sffamily\scriptsize] {$\hat B$};
  \draw[->, muted] (0,-1.85) -- (0,1.95) node[above, font=\sffamily\scriptsize] {$\hat N$};
  \draw[linhasuave] (0,0) circle (1.25);
  \draw[laranjaesc, thick]
    (0,1.25) \foreach \k in {1,...,8} { -- ({1.25*sin(\k*45)},{1.25*cos(\k*45)}) };
  \foreach \k in {0,...,7} {
    \fill[laranjaesc] ({1.25*sin(\k*45)},{1.25*cos(\k*45)}) circle (0.9pt);
  }
  \draw[verdeesc, thick, ->] (0,0) -- ({1.25*sin(45)},{1.25*cos(45)})
        node[midway, above left, font=\scriptsize] {$o$};
  \fill (0,0) circle (1.15pt) node[below left, font=\scriptsize] {$p(t)$};
  \node[font=\sffamily\scriptsize, muted] at (1.55,-1.55) {8 lados};
\end{tikzpicture}
\captionof{figure}{Seção do tubo, olhando ao longo da tangente. \(\alpha = 0\) fica sobre \(\hat N\). O octógono é o que a tecla \fn{G}, com a luz ligada, mostra facetado. A normal de cada vértice é o raio.}
\end{figura}

Oito lados é uma escolha didática. Com mais lados o tubo fica redondo também em \fn{GL\_FLAT}, e a tecla \fn{G} perde a demonstração. Com menos, o traço da logo deixa de parecer fio.

Cada lado vira uma triangle strip ao longo da volta. Para a fatia de ângulos \(\alpha_0\) e \(\alpha_1\), o laço de \(i = 0\) até \(96\) inclusive emite, em todo \(t\), o vértice de \(\alpha_0\) e o de \(\alpha_1\). O índice vai até \fn{PETAL\_SEGMENTS} para \(t\) alcançar \(2\pi\) e a curva fechar: o último ponto repete o primeiro. A última fatia fecha o octógono pelo mesmo motivo, periodicidade do seno. Há vértice repetido na costura. No modo imediato isso não tem preço.

A ordem da fita concorda com a normal radial. No primeiro triângulo, as duas arestas que saem do vértice \(\alpha = 0\) apontam aproximadamente para \(B\) e para \(T\). O produto \(B \times T\) devolve \(N\), que é exatamente \(o(0)\). Não é uma coincidência a corrigir no relatório. É a razão de a luz pegar o lado de fora do fio.

\subsection{Onde os anéis se sentam}

Os vinte anéis de uma mecha não são espaçados por igual. Um expoente maior que \(1\) em \(u^{e}\) segura os primeiros anéis perto do miolo. Sozinho, ele abre um vão grande antes da borda. A função \fn{ring\_position} devolve um pedaço desse vão no meio do caminho.

\begin{codigo}{examples/petal.c}
mid = sinf(u * MATH3D_PI);
mid = mid * mid;
return powf(u, exponent) * (1.f - 0.62f * mid) + u * (0.62f * mid);
\end{codigo}

\begin{equation}
  m(u) = \sin^{2}(\pi u),
  \qquad
  v(u) = u^{e}\,\big(1 - 0{,}62\, m(u)\big) + u\,\big(0{,}62\, m(u)\big).
  \label{eq:v}
\end{equation}
\(m\) vale \(0\) no miolo e na borda, e \(1\) no meio. Nas pontas, \(v = u^{e}\). No meio, \(62\%\) do valor vem do espaçamento linear e o resto do expoente. Para \(e = 2{,}15\) e \(u = 1/2\), \(v \approx 0{,}396\), enquanto a reta daria \(0{,}5\) e a potência pura daria \(0{,}25\).

\begin{figura}
\begin{tikzpicture}[x=12.4cm, y=7.15cm]
  \draw[linhasuave] (0,0) grid[xstep=0.25, ystep=0.25] (1,1);
  \draw[->, muted] (0,0) -- (1.05,0) node[right, font=\sffamily\scriptsize] {$u$};
  \draw[->, muted] (0,0) -- (0,1.06) node[above, font=\sffamily\scriptsize] {$v$};
  \foreach \x/\lab in {0/0, 0.25/{0{,}25}, 0.5/{0{,}5}, 0.75/{0{,}75}, 1/1}{
    \node[below, font=\scriptsize] at (\x,0) {\lab};
  }
  \foreach \y/\lab in {0.25/{0{,}25}, 0.5/{0{,}5}, 0.75/{0{,}75}, 1/1}{
    \node[left, font=\scriptsize] at (0,\y) {\lab};
  }
  \draw[muted, dashed] (0,0) -- (1,1);
  \draw[laranjaesc, thick] plot coordinates {
    (0.000,0.000) (0.050,0.002) (0.100,0.013) (0.150,0.034) (0.200,0.068)
    (0.250,0.113) (0.300,0.166) (0.350,0.225) (0.400,0.286) (0.450,0.343)
    (0.500,0.396) (0.550,0.442) (0.600,0.483) (0.650,0.521) (0.700,0.560)
    (0.750,0.604) (0.800,0.658) (0.850,0.724) (0.900,0.803) (0.950,0.896) (1.000,1.000)};
  \draw[laranja, thick] plot coordinates {
    (0.000,0.000) (0.050,0.006) (0.100,0.023) (0.150,0.051) (0.200,0.090)
    (0.250,0.138) (0.300,0.194) (0.350,0.253) (0.400,0.313) (0.450,0.370)
    (0.500,0.423) (0.550,0.471) (0.600,0.516) (0.650,0.559) (0.700,0.602)
    (0.750,0.650) (0.800,0.703) (0.850,0.765) (0.900,0.836) (0.950,0.915) (1.000,1.000)};
  \draw[verdeesc, thick] plot coordinates {
    (0.000,0.000) (0.050,0.017) (0.100,0.045) (0.150,0.083) (0.200,0.128)
    (0.250,0.179) (0.300,0.235) (0.350,0.292) (0.400,0.348) (0.450,0.403)
    (0.500,0.456) (0.550,0.506) (0.600,0.554) (0.650,0.600) (0.700,0.647)
    (0.750,0.696) (0.800,0.749) (0.850,0.806) (0.900,0.867) (0.950,0.932) (1.000,1.000)};
  \foreach \u/\v in {0.000/0.000, 0.050/0.002, 0.100/0.013, 0.150/0.034, 0.200/0.068, 0.250/0.113, 0.300/0.166, 0.350/0.225, 0.400/0.286, 0.450/0.343, 0.500/0.396, 0.550/0.442, 0.600/0.483, 0.650/0.521, 0.700/0.560, 0.750/0.604, 0.800/0.658, 0.850/0.724, 0.900/0.803, 0.950/0.896}{
    \fill[laranjaesc] (\u,\v) circle (0.55pt);
  }
  \node[font=\sffamily\scriptsize, anchor=west, laranjaesc] at (0.62,0.30) {$e = 2{,}15$};
  \node[font=\sffamily\scriptsize, anchor=west, laranja] at (0.62,0.38) {$e = 1{,}75$};
  \node[font=\sffamily\scriptsize, anchor=west, verdeesc] at (0.62,0.48) {$e = 1{,}38$};
  \node[font=\sffamily\scriptsize, anchor=west, muted] at (0.62,0.62) {linear};
\end{tikzpicture}
\captionof{figure}{Raio normalizado \(v\) contra o parâmetro \(u\). As três curvas são as três mechas. Os pontos são os vinte anéis da mecha de expoente \(2{,}15\): apertados no miolo, mais folgados na borda.}
\end{figura}

Cada mecha desloca a amostra. Com deslocamento \(\delta \in \{0,\ 1/3,\ 2/3\}\) e vinte anéis,
\[
  u_0 = (k + \delta)/20, \qquad u_1 = \min\big(1,\ (k + 1 + \delta)/20\big),
\]
e os raios usados são \(v(u_0)\) e \(v(u_1)\). Os expoentes são \(2{,}15\), \(1{,}75\) e \(1{,}38\). A mecha de expoente maior fica mais tempo perto do miolo. Os deslocamentos impedem que as três famílias caiam nas mesmas elipses. Há ainda um piso: se \(v_1 < v_0 + 0{,}006\), \(v_1\) é empurrado. Sem isso, a potência achata tanto o primeiro passo que o tubo nasce com comprimento nulo e a malha degenera.

O anel de índice \(k\) não interpola o raio ao longo da volta. Ele usa um raio só, o de \(v_0\), e gasta a volta para subir de \(z(v_0)\) até \(z(v_1)\):
\begin{equation}
\begin{aligned}
r_x(v) &= 0{,}058 + 0{,}29\, v,\\
r_y(v) &= 0{,}090 + 0{,}46\, v,\\
z(v) &= -0{,}50\,(1 - v)^{0{,}72} + \phi .
\end{aligned}
\label{eq:raio}
\end{equation}
Na borda, \(v = 1\), a elipse mede \(0{,}348\) por \(0{,}550\). No miolo, \(0{,}058\) por \(0{,}090\). A razão \(r_y/r_x\) fica por volta de \(1{,}55\) nos dois extremos: a família é alongada no eixo da pétala do começo ao fim.

\(\phi = (\delta - 0{,}5)\cdot 0{,}06\) é um empurrão minúsculo em \(z\), entre \(-0{,}03\) e \(0{,}01\). Perto da profundidade de \(0{,}50\), não separa as mechas. Inclina uma em relação à outra, e de lado os fios se cruzam em vez de coincidir.

\subsection{O funil}

\(z(0) = -0{,}50\) e \(z(1) = 0\). O miolo está mais longe da câmera; a borda, mais perto. O expoente \(0{,}72\), menor que \(1\), segura o meio da pétala para trás: em \(v = 1/2\), \(z \approx -0{,}304\), enquanto uma reta já estaria em \(-0{,}25\). A tigela é funda no miolo e vem à frente mais depressa perto da borda.

\begin{figura}
\begin{tikzpicture}[x=8.4cm, y=4.15cm, >=Stealth, line cap=round]
  \draw[->, muted] (-0.66,0) -- (0.58,0) node[right, font=\sffamily\scriptsize] {$z$};
  \draw[->, muted] (0,-0.78) -- (0,0.82) node[above, font=\sffamily\scriptsize] {$y$};
  \foreach \z/\ry in {-0.500/0.090, -0.406/0.205, -0.304/0.320, -0.184/0.435, 0.000/0.550}{
    \draw[verdeesc, line width=1.15pt] (\z,-\ry) -- (\z,\ry);
    \fill[verdeesc] (\z,0) circle[radius=1.25pt];
  }
  \draw[laranjaesc, dashed] (-0.500,0.090) .. controls (-0.34,0.18) and (-0.16,0.36) .. (0.000,0.550);
  \draw[laranjaesc, dashed] (-0.500,-0.090) .. controls (-0.34,-0.18) and (-0.16,-0.36) .. (0.000,-0.550);
  \node[font=\sffamily\scriptsize, muted, align=center] at (0.36,0.34) {câmera\\olha daqui};
  \draw[->, laranjaesc] (0.46,0) -- (0.10,0);
  \node[font=\scriptsize, above] at (-0.50,0.11) {miolo};
  \node[font=\scriptsize, above left] at (0.00,0.55) {borda};
\end{tikzpicture}
\captionof{figure}{Uma pétala no próprio espaço, vista de lado. Cada traço vertical é a altura \(2 r_y\) de um anel, no \(z\) da equação~\eqref{eq:raio}. O largo está perto da câmera. A tecla \fn{3} mostra as quatro cópias desta figura, já colocadas nas diagonais.}
\end{figura}

O número \(0{,}72\) não sai de uma fórmula da logo. É forma. A vista de lado é o lugar de julgá-lo. O mesmo vale para \(0{,}50\), a profundidade total: mais fundo, mais volume; fundo demais, e a vista frontal começa a denunciar a rampa.

\subsection{A cor não é geometria}

\begin{codigo}{examples/petal.c}
out[i] = dark[i] + (light[i] - dark[i]) * u;
\end{codigo}

\(u\) aqui é o \(v\) do anel, não o ângulo da elipse. Um tubo inteiro tem uma cor só. O degradê da logo é discreto, de tubo em tubo, do escuro no miolo ao claro na borda. As três mechas se intercalam, e o olho lê um degradê mais fino do que vinte amostras.

A cor entra por \fn{glColor3fv}, em componentes de \(0\) a \(1\). Os literais do fonte são as amostras de \(0\) a \(255\) divididas por \(255\), tiradas da referência.

\begin{center}
\captionof{table}{As quatro cores. Com a luz ligada elas viram \(k_a\) e \(k_d\), não a cor final do pixel.}
\small
\begin{tabular}{@{}l l l@{}}
\toprule
Papel & RGB em \(0..255\) & Uso \\
\midrule
laranja escuro & \(215,\ 126,\ 42\) & miolo das pétalas de cima \\
laranja claro & \(234,\ 168,\ 85\) & borda das pétalas de cima \\
verde escuro & \(96,\ 134,\ 72\) & miolo das pétalas de baixo \\
verde claro & \(150,\ 186,\ 118\) & borda das pétalas de baixo \\
\bottomrule
\end{tabular}
\end{center}

Há, portanto, duas variações em cada fio, e elas não se confundem. Ao longo da família de anéis, muda o material. Em volta de um mesmo fio, o material é constante e muda a normal. A redondez do traço é iluminação. O clarear para a borda é cor. Desligar a luz, tecla \fn{L}, conserva o degradê e apaga a redondez.

\subsection{O que cada constante muda}
\label{sec:constantes}

\begin{figura}
\captionof{table}{Constantes de forma. O valor é o do fonte. A terceira coluna é o que olhar depois de mudar.}
\small
\begin{tabularx}{\linewidth}{@{}l r X@{}}
\toprule
Constante & Valor & Se aumentar \\
\midrule
\fn{PETAL\_RINGS} & 20 & mais traços; se colarem, desça o raio do tubo \\
\fn{PETAL\_TUBE\_RADIUS} & \(0{,}0082\) & fio mais grosso \\
\fn{PETAL\_DEPTH} & \(0{,}50\) & funil mais fundo, vista \fn{3} \\
expoente de \(z\) & \(0{,}72\) & o meio da pétala vem à frente mais cedo \\
expoentes das mechas & \(2{,}15;\ 1{,}75;\ 1{,}38\) & o miolo fica mais apertado \\
\(r_x, r_y\) & eq.~\eqref{eq:raio} & lóbulo mais largo ou mais comprido \\
\fn{PETAL\_OUTWARD} & \(0{,}44\) & afasta o centro do lóbulo da origem \\
escalas \(1{,}10\) e \(0{,}84\) & & acentuam a diagonal desigual \\
fatores \(1{,}04\) e \(0{,}90\) & & afastam o lóbulo grande e aproximam o pequeno \\
\fn{PETAL\_SEGMENTS} & 96 & elipse mais lisa \\
\fn{PETAL\_SLICES} & 8 & fio mais redondo; o \fn{GL\_FLAT} fica menos legível \\
\bottomrule
\end{tabularx}
\end{figura}

\begin{voz}
Uma pétala, em uma frase: sessenta tubos elípticos, três famílias com expoentes diferentes, miolo em \(z = -0{,}5\), borda em \(z = 0\), normal de iluminação igual ao raio do tubo. \fn{main.c} só copia essa malha quatro vezes.
\end{voz}

\section{Normal, Phong e Gouraud}
\label{sec:luz}

\subsection{O modelo de Phong}

Num ponto da superfície, com normal unitária \(N\), a direção unitária \(L\) do ponto até a luz, e a direção unitária \(V\) do ponto até o observador, o modelo de Phong escreve três parcelas:
\begin{equation}
  I = I_a k_a + I_d k_d \max(N\cdot L,\ 0) + I_s k_s \max(R\cdot V,\ 0)^{n},
  \label{eq:phong}
\end{equation}
\begin{equation}
  R = 2 (N\cdot L)\, N - L .
  \label{eq:reflex}
\end{equation}
\(I_a, I_d, I_s\) são as intensidades ambiente, difusa e especular da luz. \(k_a, k_d, k_s\) são as reflectâncias do material. \(n\) é o brilho: grande, o especular fica um ponto; pequeno, fica uma mancha.

A parcela ambiente não depende da normal. Ela existe para o lado que a luz não alcança não ficar preto puro. A difusa é a lei do cosseno: superfície de frente para a luz, \(N\cdot L\) perto de \(1\), fica clara; de perfil, fica escura; de costas, o máximo com zero a apaga. A especular depende de o reflexo \(R\) apontar para o observador.

\begin{figura}
\begin{tikzpicture}[scale=1.05, >=Stealth, thick]
  \draw[linhasuave] (-2.4,0) -- (2.6,0);
  \fill[codebg] (-2.4,0) -- (2.6,0) -- (2.6,-0.55) -- (-2.4,-0.55) -- cycle;
  \draw[muted] (-2.4,0) -- (2.6,0);
  \coordinate (P) at (0,0);
  \draw[verdeesc, ->] (P) -- (0,1.7) node[above, font=\scriptsize] {$N$};
  \draw[laranjaesc, ->] (P) -- (1.35,1.15) node[right, font=\scriptsize] {$L$};
  \draw[laranja, ->] (P) -- (-0.55,1.45) node[left, font=\scriptsize] {$R$};
  \draw[ink, ->] (P) -- (-1.25,0.85) node[left, font=\scriptsize] {$V$};
  \draw[verdeesc, dashed, ->] (P) -- (0.22,1.42);
  \node[font=\scriptsize, verdeesc, anchor=west] at (0.32,1.05) {$H$};
  \fill (P) circle (1.3pt);
\end{tikzpicture}
\captionof{figure}{\(R\) é o reflexo de \(L\) em torno de \(N\). \(H\) é o meio do caminho entre \(L\) e \(V\), que é o vetor que a OpenGL usa no lugar de \(R\).}
\end{figura}

\subsection{O que a OpenGL calcula}

O pipeline fixo não avalia \(R\cdot V\). Ele avalia o modelo de Blinn--Phong. O vetor do meio é
\begin{equation}
  H = \frac{L + V}{\|L + V\|},
  \qquad
  \text{especular} = \max(N\cdot H,\ 0)^{n}.
  \label{eq:blinn}
\end{equation}
\(H\) aponta para onde a normal precisa apontar para que a luz rebata no observador. Quando \(N\) coincide com \(H\), o especular é máximo. O enunciado aceita a rotina da OpenGL. No relatório, a fórmula escrita é a~\eqref{eq:phong}, com uma frase dizendo que a implementação usa o halfway.

Há três guardas que a fórmula de quadro não mostra e a rotina aplica. Se \(N\cdot L < 0\), a luz está do outro lado da superfície e o especular também vai a zero, não só a difusa. A atenuação com a distância não foi configurada: os coeficientes ficam no padrão, constante \(1\) e linear e quadrático \(0\), então a lâmpada não enfraquece com a distância. O que muda o brilho de um fio para o outro é o ângulo, não o fato de uma pétala estar mais longe da lâmpada. E \fn{GL\_LIGHT\_MODEL\_LOCAL\_VIEWER} ligado faz \(V\) ir do vértice até a câmera, em vez de usar a direção fixa \((0, 0, 1)\).

\fn{GL\_SHININESS} é o \(n\). O valor no fonte é \(48\), dentro da faixa \(0\) a \(128\) do pipeline fixo.

\begin{codigo}{examples/main.c · os coeficientes}
glLightModelfv(GL_LIGHT_MODEL_AMBIENT, global_ambient); /* 0.46, 0.46, 0.44 */
glLightfv(GL_LIGHT0, GL_AMBIENT, ambient);              /* 0.22 */
glLightfv(GL_LIGHT0, GL_DIFFUSE, diffuse);              /* 0.88 */
glLightfv(GL_LIGHT0, GL_SPECULAR, specular);            /* 0.28 */
glMaterialfv(GL_FRONT_AND_BACK, GL_SPECULAR, mat_specular); /* 0.22 */
glMaterialf(GL_FRONT_AND_BACK, GL_SHININESS, 48.f);
glLightModeli(GL_LIGHT_MODEL_COLOR_CONTROL, GL_SEPARATE_SPECULAR_COLOR);
\end{codigo}

\(k_a\) e \(k_d\) não estão nessa lista. Eles vêm da cor da pétala, porque \fn{GL\_COLOR\_MATERIAL} está ligado em \fn{GL\_AMBIENT\_AND\_DIFFUSE}. \(k_s\) é o cinza \(0{,}22\) do material, igual nas quatro pétalas. A luz especular é fraca de propósito, \(0{,}28\): um fio fino com especular branco forte vira um risco de giz e a cor da logo some.

\fn{GL\_SEPARATE\_SPECULAR\_COLOR} interpola o especular à parte e soma depois da cor difusa. Sem isso, o especular é somado à cor ainda no vértice e o resultado é grampeado em \(1\) antes de viajar pelo triângulo. Num laranja que já está alto, o brilho não tem onde aparecer e a parcela de Phong fica invisível na apresentação. O controle não transforma Gouraud em Phong por pixel. O brilho continua calculado nos vértices.

\fn{GL\_LIGHT\_MODEL\_TWO\_SIDE} liga a equação também na face de trás, invertendo a normal. O material de trás é o mesmo. A fita já nasce coerente com a normal; o estado de dois lados cobre o que ainda assim aparecer virado, em especial na silhueta do tubo.

A posição da luz é reenviada todo quadro. Guardá-la uma vez no \fn{init}, com a modelview ainda identidade, grava a lâmpada em coordenada de olho sem a translação da câmera. O ponto \((1{,}8,\ 2{,}4,\ 5)\) ficaria atrás do observador, porque \(z\) positivo é o lado de trás da câmera. O fonte repete a chamada no lugar certo e essa armadilha fica desarmada.

\subsection{Gouraud é outra decisão}

O modelo de Phong diz como \(I\) sai de \(N\), \(L\) e \(V\). A tonalização diz onde essa conta é feita, e o que se interpola no meio do triângulo.

\begin{center}
\captionof{table}{Três nomes que a arguição troca, e só um deles está na tecla \fn{G}.}
\small
\begin{tabularx}{\linewidth}{@{}l X X@{}}
\toprule
Nome & Onde a conta de luz acontece & O que anda pelo triângulo \\
\midrule
Gouraud, \fn{GL\_SMOOTH} & em cada vértice & a cor já calculada \\
Constante, \fn{GL\_FLAT} & num vértice da faceta & uma cor só \\
Phong por pixel & em cada fragmento & a normal, e a luz é refeita \\
\bottomrule
\end{tabularx}
\end{center}

A tecla \fn{G} alterna as duas primeiras. A terceira, Phong shading, o enunciado não pede e o pipeline fixo não faz. Phong, nesta entrega, é o modelo. Gouraud é o transporte da cor.

Com \fn{GL\_FLAT}, cada triângulo da fita usa uma normal só, a do vértice provocador. Os oito lados do tubo aparecem como facetas. Com \fn{GL\_SMOOTH}, as cores dos vértices --- já iluminadas com normais diferentes --- são interpoladas, e o octógono parece redondo. A diferença só aparece com a luz ligada. Com a luz desligada, todos os vértices de um anel têm o mesmo \fn{glColor}, e interpolar uma constante não muda nada. Quem aperta \fn{G} no escuro e não vê efeito não quebrou a tecla.

O aramado, tecla \fn{W}, é \fn{glPolygonMode(GL\_FRONT\_AND\_BACK, GL\_LINE)}. Ele mostra a fita e a seção mesmo sem luz. É o jeito de confirmar que o traço tem oito lados e não é uma fita chata.

\begin{voz}
Frase pronta: a luz avalia Phong--Blinn no vértice; \fn{GL\_SMOOTH} interpola essa cor; \fn{GL\_FLAT} repete a cor da faceta; sombrear Phong de verdade interpolaria a normal e refaria a conta no pixel, e não é o que a tecla \fn{G} faz.
\end{voz}

\section{O Z-buffer}
\label{sec:z}

Cada pixel do framebuffer guarda, além da cor, uma profundidade. O teste desta cena é o mais simples:

\begin{enumerate}
\item A limpeza enche o buffer com \(1\), o longe.
\item Para cada fragmento, compara-se a profundidade dele com a armazenada.
\item Se a nova for menor (\fn{GL\_LESS}), a cor substitui a antiga e a profundidade também.
\item Se não for, o fragmento é descartado. A cor que já estava lá permanece.
\end{enumerate}

Isso é \fn{glEnable(GL\_DEPTH\_TEST)}, \fn{glDepthFunc(GL\_LESS)} e o bit \fn{GL\_DEPTH\_BUFFER\_BIT} na limpeza. A tecla \fn{Z} só liga e desliga o teste. A função de comparação continua \fn{GL\_LESS}. Profundidades iguais não passam: o primeiro fragmento permanece.

A ordem de desenho é a do laço: pétala \(0\), depois \(1\), \(2\) e \(3\). Com o teste ligado, a ordem quase não importa. Ganha quem está mais perto da câmera. Com o teste desligado, ganha quem foi desenhado por último, mesmo que esteja atrás. Onde as malhas se cruzam --- e elas se cruzam, porque a ponta interna atravessa a origem --- a diferença fica óbvia. Na vista frontal o efeito existe, mas é mais frio, porque um lóbulo não tapa o outro inteiro. A foto que convence é a vista inclinada, lado a lado, \fn{Z} ligado e \fn{Z} desligado, sem mexer no ângulo entre as duas.

\fn{glPushMatrix} não é visibilidade. Back-face culling, \fn{GL\_CULL\_FACE}, também não entra: com uma normal inconsistente ele apaga um pedaço de tubo e atrapalha mais do que ajuda. O item é o buffer de profundidade.

Dois esquecimentos deixam o teste ``ligado'' e a imagem errada. Sem \fn{GLUT\_DEPTH} na criação da janela, não há buffer. Sem o bit de profundidade no \fn{glClear}, o quadro anterior continua morando no buffer, e a cena em movimento mancha ou fura. Limpar a profundidade com \(0\), em vez de \(1\), também esvazia a imagem: nenhum fragmento tem profundidade menor que zero, e \fn{GL\_LESS} rejeita todos.

\begin{voz}
O Z-buffer guarda quem está na frente. Ele não ordena os objetos antes de desenhar, não conhece a pétala e não substitui a matriz. Compara um número por pixel, depois que a projeção já transformou \(z\).
\end{voz}

\section{A demonstração, tecla por tecla}
\label{sec:teclas}

A janela abre em yaw \(-32^\circ\) e pitch \(26^\circ\), distância \(5{,}2\), perspectiva, luz ligada, Gouraud, Z-buffer ligado, escala de pétala \(1\). Essa é a vista em que o volume aparece. A ordem abaixo cabe em cinco minutos e percorre os itens do trabalho.

\begin{enumerate}
\item \fn{F}. Quatro lóbulos, laranja em cima, verde embaixo, a diagonal principal maior. Modelagem.
\item Setas, ou o arrasto com o botão esquerdo. Depois \fn{W}. A seção do tubo aparece. \fn{W} de novo volta ao preenchido.
\item \fn{+} duas vezes e \fn{M} no terminal, com a conta da seção das pétalas ao lado. A escala mudou o ponto; a translação da matriz impressa, não.
\item \fn{P} e \fn{O} sem tocar nas setas. Ortográfica um pouco menor, paralelas conservadas.
\item \fn{L} apaga a redondez e deixa o degradê. \fn{L} de novo. \fn{G} facetá o fio. Uma frase: Phong é a conta, Gouraud é a interpolação.
\item Espaço. O brilho caminha porque a luz foi declarada antes da rotação.
\item Inclinar de novo e alternar \fn{Z} no mesmo enquadramento.
\end{enumerate}

\begin{center}
\captionof{table}{Todas as entradas que \fn{keyboard}, \fn{special\_keys} e o mouse reconhecem.}
\small
\begin{tabularx}{\linewidth}{@{}l X@{}}
\toprule
Entrada & Efeito \\
\midrule
arrasto, setas & yaw e pitch; seta soma \(4^\circ\), arrasto soma \(0{,}35^\circ\) por pixel \\
roda, \fn{[} \fn{]}, Page Up/Down & distância da câmera; aumentar a distância afasta \\
\fn{1}, \fn{F} & frontal, yaw \(0\), pitch \(0\), conserva o zoom \\
\fn{2} & a vista inicial, \(-32^\circ\) e \(26^\circ\) \\
\fn{3} & lado, yaw \(90^\circ\), pitch \(8^\circ\): o funil \\
\fn{4} & de cima, pitch \(80^\circ\) \\
\fn{R} & volta à vista inicial e à distância \(5{,}2\) \\
\fn{P}, \fn{O} & perspectiva, ortográfica \\
\fn{W} & aramado ou preenchido \\
\fn{G} & \fn{GL\_SMOOTH} ou \fn{GL\_FLAT} \\
\fn{L} & luz ligada ou cor direta \\
\fn{Z} & teste de profundidade \\
\fn{+}, \fn{-} & escala da pétala, passo \(0{,}05\), presa a \([0{,}3,\ 2{,}2]\) \\
espaço & gira o yaw sozinho; para ao escolher uma vista \\
\fn{M} & imprime \(S\), \(T\), \(R_z\) e o ponto \((0{,}20,\ 0,\ 0)\) \\
Esc & encerra \\
\bottomrule
\end{tabularx}
\end{center}

O arrasto soma \(y\) de tela ao pitch. No GLUT, \(y\) cresce para baixo, então arrastar para baixo aumenta o pitch e revela a parte de cima, coerente com a tecla \fn{4}. A rotação automática não usa relógio: cada chamada ociosa soma \(0{,}35^\circ\) ao yaw. A velocidade depende de quantas vezes o GLUT entrega o ócio, não de um passo de tempo fixo.

\section{Erros que custam uma tarde}
\label{sec:erros}

\begin{enumerate}
\item Matriz preenchida por linha e passada a \fn{glMultMatrixf}. A translação cai nos índices \(3\), \(7\) e \(11\).
\item Ordem dos \fn{apply\_matrix} invertida, e o ângulo alterado até ``parecer'' certo. A tecla \fn{M} desmente.
\item \fn{glMatrixMode} deixado em \fn{GL\_PROJECTION}. A câmera é apagada no quadro seguinte pela identidade da projeção.
\item \fn{glRotatef} ou \fn{gluPerspective} na pressa. A tela acerta e o item zera.
\item \fn{glLoadIdentity} ausente. As transformações acumulam e a logo sai do quadro em poucos segundos.
\item Pop faltando. A segunda pétala herda a rotação da primeira.
\item Luz configurada só no \fn{init}. Ela gruda na câmera, e no lugar errado.
\item \fn{glColor} com iluminação ligada e \fn{GL\_COLOR\_MATERIAL} desligado. A logo fica cinza.
\item Normal ausente. Com a luz ligada o tubo fica chapado ou preto. O aramado ainda mostra a seção; a luz, não.
\item Escala alta na ordem certa. Os lóbulos crescem para o centro e se atravessam. Não é bug da rotação.
\item \fn{GLUT\_DEPTH} ou o bit de limpeza da profundidade esquecidos. A tecla \fn{Z} mente.
\item Parâmetro chamado \fn{near} ou \fn{far} num arquivo que inclui \fn{windows.h}.
\item Julgar a tecla \fn{G} com a luz desligada. Não há o que interpolar.
\end{enumerate}

\section{Exercícios}
\label{sec:ex}

Façam a conta antes de olhar a conferência. Os números que dependem da tecla \fn{M} devem bater com o terminal de outra pessoa, não só com esta página.

\begin{enumerate}
\item Em que índice de \fn{float m[16]} fica \(t_y\)? E \(s_y\)? E o \(-1\) da perspectiva?
\item Com a escala da tecla \fn{M} em \(1\), sem ter apertado \fn{+}, qual é o ponto \((0{,}20,\ 0,\ 0)\) depois de \(R_z(45^\circ)\, T\, S\), em três casas?
\item Apertar \fn{+} uma vez leva a escala a \(1{,}05\). Calcule o novo ponto antes de apertar \fn{M}. A última coluna de \(R_z T\) muda?
\item Por que a pétala de \(+45^\circ\) não coincide com a matriz que \fn{M} imprimiu, mesmo com a escala global em \(1\)?
\item Troquem, no papel, a ordem para \(S\, T\, R_z\). Onde fica o centro das quatro pétalas? O que a tecla \fn{+} passa a fazer com a distância ao centro?
\item \(\varphi\) cai de \(40^\circ\) para \(20^\circ\) e \(d\) continua \(5{,}2\). A logo cresce ou encolhe? A ponta, a cerca de \(0{,}75\) em \(Y\), ainda cabe na meia-altura visível?
\item A escala da pétala é uniforme. Por que \fn{GL\_NORMALIZE} ainda muda a luz?
\item A luz é declarada depois de \(R_y\). O que o brilho faz quando o espaço está ligado?
\item Com o Z-buffer desligado, dois fragmentos caem no mesmo pixel. Qual fica?
\item Calcule \(v(1/2)\) para o expoente \(2{,}15\) pela equação~\eqref{eq:v}. Por que esse anel não está na metade do raio?
\item A normal da seção, no ponto direito da elipse, aponta para o centro. A normal de \fn{glNormal3f} aponta para onde?
\item Qual índice \(i\), qual ângulo e qual cor formam o lóbulo superior esquerdo?
\end{enumerate}

\newpage
\subsection*{Conferência}

\begin{enumerate}
\item \(t_y\) fica em \(m[13]\). \(s_y\) fica em \(m[5]\). O \(-1\) da perspectiva fica em \(m[11]\), última linha e terceira coluna.

\item \((-0{,}170,\ 0{,}453,\ 0{,}000)\). A conta crua dá \((-0{,}1697,\ 0{,}4525,\ 0)\); o \fn{printf} arredonda a três casas.

\item \(S p = (0{,}210,\ 0,\ 0)\), depois da translação \((0{,}210,\ 0{,}440,\ 0)\).
\(x' = 0{,}707107\,(0{,}210 - 0{,}440) = -0{,}1626\),
\(y' = 0{,}707107\,(0{,}210 + 0{,}440) = 0{,}4596\).
Três casas: \((-0{,}163,\ 0{,}460,\ 0{,}000)\). A coluna da translação não muda, porque \(T\) continua \((0,\ 0{,}44,\ 0)\) e \(R_z\) é o mesmo. \(S\) só alterou o ponto, que não está na origem.

\item A pétala de \(+45^\circ\) é a grande: escala \(1{,}10\) por cima da escala global, distância \(0{,}44 \times 1{,}04 = 0{,}4576\). A tecla \fn{M} usa escala global e distância \(0{,}44\). O ângulo coincide; o tamanho e o centro, não.

\item \(R_z\) gira a pétala em torno da origem e o centro continua lá. \(T\) manda as quatro para \((0,\ d,\ 0)\). \(S\), por fora, escala em torno da origem do mundo, então a distância \(d\) cresce junto com o lóbulo. As quatro empilham no \(+Y\).

\item Cresce. A meia-altura visível passa de \(5{,}2\tan 20^\circ \approx 1{,}89\) para \(5{,}2\tan 10^\circ \approx 0{,}92\). A ponta em \(0{,}75\) ainda cabe.

\item A inversa transposta de uma escala uniforme \(s\) multiplica a normal por \(1/s\). A direção fica certa e o comprimento deixa de ser \(1\), então \(N\cdot L\) e o expoente especular saem errados. \fn{normalize3} já tinha passado, lá no objeto, antes da modelview.

\item A rotação do objeto entra na lâmpada. O brilho gruda no mesmo ponto da pétala e gira com ela, em vez de caminhar pela superfície.

\item O que foi desenhado por último. No laço, o último é \(i = 3\), ângulo \(225^\circ\), verde, lóbulo grande, inferior direito. Isso vale nos pixels em que essa malha realmente escreveu. Onde ela não cobre, ganha o anterior, e assim por diante.

\item \(m(1/2) = 1\), então \(v = (1/2)^{2{,}15}\cdot 0{,}38 + (1/2)\cdot 0{,}62\).
\((1/2)^{2{,}15} \approx 0{,}226\), e \(v \approx 0{,}086 + 0{,}310 = 0{,}396\).
A potência puxa para o miolo; o seno devolve parte da folga, não toda.

\item Para fora do fio. No ponto direito, \(o(0) = \hat N\) aponta para o centro da elipse, que é o lado interno do tubo. \(o(\pi) = -\hat N\) aponta para o lado externo. Os dois são normais de superfície, cada uma no seu \(\alpha\).

\item Índice \(1\), ângulo \(+45^\circ\), laranja, escala \(1{,}10\), distância \(0{,}4576\). O centro cai em \(x < 0\) e \(y > 0\).
\end{enumerate}

\section{Mapa do código}

\begin{center}
\captionof{table}{Onde cada assunto mora. A função é o que se cita na arguição; o arquivo é o que se abre.}
\small
\begin{tabularx}{\linewidth}{@{}l l X@{}}
\toprule
Assunto & Função & Arquivo \\
\midrule
Janela, quadro, teclas & \fn{main}, \fn{display} & \fn{main.c} \\
Luz e material & \fn{setup\_lighting} & \fn{main.c} \\
Quatro cópias & \fn{draw\_petals} & \fn{main.c} \\
Conta da tecla \fn{M} & \fn{print\_matrices} & \fn{main.c} \\
Produto e layout & \fn{mat4\_multiply} & \fn{math3d.c} \\
Afins & \fn{mat4\_make\_*} & \fn{math3d.c} \\
Projeções & \fn{mat4\_make\_perspective}, \fn{mat4\_make\_ortho} & \fn{math3d.c} \\
Porta da OpenGL & \fn{apply\_matrix} & \fn{math3d.c} \\
Curva e triedro & \fn{ellipse\_frame} & \fn{petal.c} \\
Fio & \fn{tube\_vertex}, \fn{draw\_ring} & \fn{petal.c} \\
Famílias & \fn{ring\_position}, \fn{draw\_strand} & \fn{petal.c} \\
Cabeçalho Windows & macros \fn{near}/\fn{far} & \fn{gl\_platform.h} \\
\bottomrule
\end{tabularx}
\end{center}

O que cada pessoa precisa conseguir dizer sem abrir o plano:

\begin{itemize}
\item a ordem \(R_z T S\), e por que a última chamada age primeiro;
\item o índice \(m[c\cdot 4 + r]\), com a translação em \(12\), \(13\), \(14\);
\item a diferença entre \(w = -z\) e \(w = 1\);
\item a cadeia \(T_{\mathrm{cam}} R_x R_y R_z T S\), e o lugar da \fn{glLightfv} dentro dela;
\item o tubo: elipse, triedro, \(\,v = p + \rho\, o(\alpha)\), normal igual a \(o\);
\item Phong no papel, Blinn na OpenGL, Gouraud na tecla \fn{G};
\item o teste \fn{GL\_LESS} contra a ordem de desenho.
\end{itemize}

\section{Glossário}

\begin{description}[
  leftmargin=0pt,
  style=nextline,
  font=\sffamily\bfseries,
  itemsep=0.35em]
\item[Aspect]
Largura da janela dividida pela altura. Entra em \(f_{\!v}/a\) na perspectiva e na largura do retângulo ortográfico. Sem ele, a elipse estica quando a janela muda de forma.

\item[Blinn--Phong]
O modelo de especular que a OpenGL clássica avalia: \((N\cdot H)^{n}\), com \(H\) unitário na direção de \(L + V\).

\item[Column-major]
A matriz \(4\times 4\) armazenada coluna por coluna. É o layout de \fn{glMultMatrixf} e de \fn{m[16]} neste programa.

\item[Coordenada homogênea]
O ponto \((x, y, z)\) escrito \((x, y, z, 1)\), de modo que a translação caiba numa matriz.

\item[Divisão perspectiva]
Trocar \((x, y, z, w)\) por \((x/w,\ y/w,\ z/w)\). É o passo em que o longe fica menor.

\item[Gouraud]
Iluminação nos vértices e interpolação da cor dentro do triângulo. No fonte, \fn{GL\_SMOOTH}.

\item[Modelview]
A matriz que leva o vértice do objeto ao olho. Aqui ela acumula câmera, yaw, pitch e a cópia da pétala.

\item[Modo imediato]
Emitir vértices com \fn{glBegin} e \fn{glEnd}, em vez de armazená-los num buffer.

\item[NDC]
Normalized device coordinates. O cubo \([-1, 1]^3\) depois da divisão por \(w\) e antes do viewport.

\item[Normal]
Vetor perpendicular à superfície. No tubo, é o raio da seção, de comprimento \(1\). A OpenGL transforma esse vetor pela inversa transposta da modelview.

\item[Pipeline fixo]
A OpenGL em que luz, matriz e tonalização são estado da máquina, não shaders.

\item[Projeção]
A matriz \(P\) que leva o olho ao recorte. Perspectiva ou ortográfica. Não move a câmera.

\item[Triangle strip]
Primitiva em que cada vértice novo fecha um triângulo com os dois anteriores. É como cada fatia do tubo viaja ao longo da elipse.

\item[Viewport]
O retângulo de pixels que recebe o quadrado \([-1, 1]\) do NDC. No programa, a janela inteira.

\item[Z-buffer]
Memória de profundidade por pixel. Com \fn{GL\_LESS}, fica quem está mais perto da câmera.
\end{description}

\vspace{1.2em}
{\color{linhasuave}\rule{\textwidth}{0.6pt}}
\vspace{0.6em}
{\sffamily\footnotesize\color{muted}
Os números desta apostila foram lidos de \fn{examples/main.c}, \fn{examples/math3d.c} e \fn{examples/petal.c}. Constantes de forma sem nome no fonte --- os coeficientes \(0{,}058\), \(0{,}29\), \(0{,}62\), \(0{,}72\) --- estão nas expressões de \fn{draw\_strand} e \fn{ring\_position}. Se algum desses literais mudar, a conta da seção correspondente muda com ele; a álgebra, não.
\par}

\end{document}
